\chapter{Fixed-Point, State, and Protocol Contract}
\label{app:fixed-point}

This appendix collects the bit-level constants required to reproduce the executed
controller without inferring them from prose.

\section{State and Policy Storage}

\begin{table}[H]
\centering
\caption{Adaptive-controller storage layout.}
\label{tab:rom-layout}
\begin{tabularx}{\textwidth}{@{}p{0.24\textwidth}rrX@{}}
\toprule
\textbf{Storage} & \textbf{Depth} & \textbf{Useful width} & \textbf{Ordering/use}\\
\midrule
Policy ROM & 256 & 2 bits & State IDs 0--255; deployed H4 actions\\
Profile ROM & 4 & 18 bits & Profiles $\Pi_0$--$\Pi_3$; deployed payloads\\
Training Q-table & 1,024 & 16 bits & State-major/action-minor signed Q8.7; not deployed\\
\bottomrule
\end{tabularx}
\end{table}

The radix-4 address is
\begin{equation}
z=(((b_K\times4)+b_Q)\times4+b_L)\times4+b_I.
\end{equation}
The two-bit fields occupy $[7{:}6]$, $[5{:}4]$, $[3{:}2]$, and $[1{:}0]$ in that
order. The deployed policy contains action counts $(145,43,44,24)$ for
$(\Pi_0,\Pi_1,\Pi_2,\Pi_3)$.

\section{Threshold Edge Contract}

\begin{table}[H]
\centering
\small
\caption{Real and encoded threshold constants.}
\label{tab:appendix-thresholds}
\begin{tabularx}{\textwidth}{@{}p{0.25\textwidth}cccX@{}}
\toprule
\textbf{Feature} & \textbf{Threshold 1} & \textbf{Threshold 2} &
\textbf{Threshold 3} & \textbf{UNORM16 codes}\\
\midrule
Minimum key occupancy & 0.25 & 0.50 & 0.75 & 16,384; 32,768; 49,151\\
Bottleneck route quality & 0.90 & 0.93 & 0.96 & 58,982; 60,948; 62,914\\
Offered request load & 0.25 & 0.50 & 0.75 & 16,384; 32,768; 49,151\\
Utilization imbalance & 0.125 & 0.25 & 0.50 & 8,192; 16,384; 32,768\\
\bottomrule
\end{tabularx}
\end{table}

The four-bin function is left-closed at every threshold: equality enters the higher
bin. Testing below, at, and above every one of the 12 thresholds gives 36 directed
threshold-edge cases.

\section{Numeric Formats and Special Codes}

\begin{table}[H]
\centering
\small
\caption{Fixed-point formats used by the physical candidate shell.}
\label{tab:appendix-numeric-formats}
\begin{tabularx}{\textwidth}{@{}p{0.25\textwidth}p{0.19\textwidth}p{0.20\textwidth}X@{}}
\toprule
\textbf{Quantity} & \textbf{Format} & \textbf{Scale/range} & \textbf{Rule}\\
\midrule
State and route quality & 16-bit UNORM16 & $0$--$65535\leftrightarrow0$--$1$ &
Round half up and clip\\
Path cost & 32-bit UQ16.16 & Scale $2^{16}$ & Maximum finite
\texttt{FFFFFFFE}; infinity \texttt{FFFFFFFF}\\
Candidate utilization & 32-bit UQ16.16 & Scale $2^{16}$ & Unsigned saturation\\
Hop count & 3-bit unsigned & 0--7 & Integer edge count\\
Candidate slot & 2-bit unsigned & 0--3 & Stable final tie-break\\
Tchebycheff ratio & 2-bit unsigned & 0--3 & Common scaling omitted\\
Weighted score & 18-bit unsigned & 0--262143 & Maximum of three products\\
Key count & 16-bit unsigned integer & Trace capacity 16 & Synthetic environment\\
Synthetic key rate & 16-bit UQ8.8 & Scale $2^8$ & Software/export contract\\
$\alpha$ coefficient & 3-bit U2.1 & Half-unit increments & Stored in profile word\\
\bottomrule
\end{tabularx}
\end{table}

For nonconstant range $m<M$, normalization is
\begin{equation}
N(x)=\operatorname{round}_{\rm half\ up}
\left(\frac{(x-m)65535}{(M-m)+1}\right).
\end{equation}
Equal ranges map to zero. The comparator ranking is
$(\mathrm{score},\mathrm{path\ cost},\mathrm{hops},\mathrm{slot})$, all ascending.

\section{Profile Payloads}

The 18-bit packing order is
\[
\alpha_1[2{:}0]\,\alpha_2[2{:}0]\,\alpha_3[2{:}0]\,\alpha_4[2{:}0]\,
\rho_c[1{:}0],\rho_q[1{:}0],\rho_u[1{:}0].
\]
The payloads are \texttt{13329}, \texttt{1B516}, \texttt{14399}, and
\texttt{0A2A5}. In the physical candidate evaluator, only the lower six ratio bits
affect selection because path cost is supplied precomputed. The alpha fields remain
part of the registered semantic/export word.

\section{H3 Priority Function}

For completeness, the exact first-match function is
\begin{equation}
a_{H3}(z)=\begin{cases}
1,&b_K=0,\\
2,&b_K\ne0\land b_Q\le1,\\
1,&b_K\ne0\land b_Q>1\land b_I\ge2,\\
3,&b_L=3\land b_K\ge2\land b_Q\ge2\land b_I\le1,\\
0,&\text{otherwise}.
\end{cases}
\end{equation}
The order is significant; the mathematical guards make the priority explicit.

\section{Synchronous and UART Protocol}

A request is accepted only on \texttt{start \&\& ready}. Start while busy is
ignored and does not alter the active request. Reset clears busy/done, cycle count,
winner state, profile history, and dwell state. Done is a one-cycle pulse with all
result fields stable.

\begin{table}[H]
\centering
\caption{Status and failure-result contract.}
\label{tab:appendix-status-contract}
\begin{tabularx}{\textwidth}{@{}c p{0.25\textwidth}X@{}}
\toprule
\textbf{Code} & \textbf{Name} & \textbf{Returned fields}\\
\midrule
0 & \texttt{OK} & Selected slot 0--3, computed score and route-quality code\\
1 & \texttt{NO\_PATH} & Selected path 255; score and route quality zero; valid
controller state/profile retained\\
2 & \texttt{INVALID\_REQUEST} & Selected path 255; score and route quality zero;
state/profile zero\\
\bottomrule
\end{tabularx}
\end{table}

The 44-byte record is
\begin{center}
\texttt{QF6R,1,TTTT,P,SSS,R,PPP,CCCCCC,FFFFF,LLL,E\textbackslash r\textbackslash n}.
\end{center}
The nine numeric fields are test ID, policy ID, state ID, profile ID, selected path,
score, route-quality code (registered field name bottleneck fidelity), cycles, and
status. Production UART is 100~MHz/115,200 baud, 8N1, idle high, no flow control.
Bytes are emitted left to right; each byte sends least-significant data bit first.

\section{Latency Boundary and Observed Ranges}

Cycle count is $N_{\rm done}-N_{\rm accept}$. H0 normal results are exactly two
cycles. H1 normal results span 7--309 cycles, H2/H3 span 8--310, and H4 spans
10--312. Mean values are 2.000, 201.763, 202.763, 202.763, and 204.763 cycles.
Host parsing, UART serialization, operating-system scheduling, JTAG, programming,
KMS access, candidate generation, and route installation are excluded.
